\section{Cell and angular information from the original spectral gap}
\label{sec:peps_angular_information_from_gap}

\begin{theorem}[Cell information from the original spectral gap]
    \label{thm:peps_cell_conditional_information_from_gap}
    \lean{TNLean.PEPS.Approximation.exists_cell_conditionalMutualInformation_le_square}
    \leanok
    \uses{def:polynomial_peps_model,def:area_law_ground_entropy}
    Fix $q\in\mathbb N$ with $q\ge1$, $J\ge0$, $\Delta>0$ and $c>0$.
    There are $C_1>0$ and $K\ge0$ with the following property.
    Let $s\ge1$, let $o\in\mathbb R^2$, and put
    $Q_i=o+si+[0,s)^2$ for $i\in\mathbb Z^2$.
    Fix $i_0\in\mathbb Z^2$ and a set $A\subseteq\mathbb Z^2$ with $i_0\notin A$.
    Let $H$ be a Hamiltonian of the square-grid model on
    $\Lambda_L=\{1,\ldots,L\}^2$, with interaction bound $J$, and let $\Omega$
    be its unit ground vector with full-system gap $\Delta$.
    Set $G=\Lambda_L\cap Q_{i_0}$ and
    $B=\Lambda_L\cap\bigcup_{i\in A}Q_i$.
    For every $E\subseteq\Lambda_L\setminus(G\cup B)$ satisfying
    $\operatorname{dist}_\infty(\overline{Q_{i_0}},E)\ge cs$, one has
    \begin{align}
        I_\Omega(G:E\mid B) & \le C_1(1+\log s)^K.
        \notag
    \end{align}
    The constants precede the scale, grid origin, conditioning cells,
    system size, Hamiltonian, ground vector and exterior region.
    This is the cell estimate of
    \cite[Lemma~3.2]{OpenAI2026PolynomialPEPS}.
\end{theorem}
\begin{proof}\leanok
    \uses{thm:area_law_sublinear_collar,
        thm:peps_square_cell_conditional_information_of_collar}
    Apply the sublinear collar theorem with range $R=1$, clearance $D_0=13$
    and template constant $C_{\rm tpl}=24$. The required inequalities are
    $2R+10=12<13$ and $24\le C_{\rm tpl}$.
    The conditional square-cell estimate then applies to the same ground
    vector, cell, conditioning region and exterior.
\end{proof}

\begin{theorem}[Information under angular separation]
    \label{thm:peps_angular_information_from_gap}
    \lean{TNLean.PEPS.Approximation.angularInformationBound}
    \leanok
    \uses{def:polynomial_peps_model,def:peps_angular_separation}
    Fix $q\in\mathbb N$ with $q\ge1$, $J\ge0$ and $\Delta>0$.
    For every $a>0$, $A_0\ge1$, $m_0\in\mathbb N$ and $p>0$, there are
    $D,L_0\in\mathbb R$ with the following property.
    Let $L\ge L_0$ be a natural number, let $H$ be a Hamiltonian of the
    square-grid model on $\Lambda_L=\{1,\ldots,L\}^2$ with interaction bound $J$,
    and let $\Omega$ be its unit ground vector with full-system gap $\Delta$.
    Let $\mathcal V\subseteq\mathbb R^2$ consist of at most $m_0$ marked points,
    let $X,E\subseteq\Lambda_L$ be disjoint, and let $n\in\mathbb R$ satisfy
    $D\log L\le n\le A_0L$ and $\operatorname{diam}_\infty X\le A_0n$.
    Put $d_{\mathcal V}(x)=\min_{v\in\mathcal V}|x-v|_\infty$, with
    $d_{\varnothing}=+\infty$.
    If every $x\in X$ and $y\in E$ satisfy
    \begin{align}
        |x-y|_\infty & \ge a\max\{D\log L,\min(n,d_{\mathcal V}(x))\},
        \notag
    \end{align}
    then $I_\Omega(X:E)\le L^{-p}$.
    The constants are chosen before the system size, Hamiltonian, ground
    vector, marked points, physical regions and scale.
    This is~\cite[Proposition~3.4]{OpenAI2026PolynomialPEPS}.
\end{theorem}
\begin{proof}\leanok
    \uses{thm:area_law_sublinear_collar,thm:peps_angular_information_of_collar}
    The sublinear collar theorem with $R=1$, $D_0=13$ and $C_{\rm tpl}=24$
    follows from the original model and gap assumptions. Apply the conditional
    angular information theorem with these fixed parameters.
    Its constants have precisely the stated order of quantifiers.
\end{proof}
